\begin{afterword}
\summaryhead

The post asks what would be missing if religion disappeared, and answers that the largest gap would be community, not ideals or morals. The author guesses that many people who stay religious without quite believing stay for their neighbours, while granting that this is something the author may not understand well.

Churches and friendly workplaces provide community only as a side effect; neither was designed for it. Looked at as a community, a church might do better without early Sunday mornings, formal clothes and weekly sermons on one theme, and with buildings shared between groups, more help for newcomers, and experiments to find what works. A community needs something in common, and a shared goal is the natural choice.

The proposal is that rationalist communities form as local ``taskforces'' around work on the world's problems, each built around the most specific interest that can support ``a decent-sized band.'' Adult learning is another possible center. The author calls the whole idea ``Maybe. Probably.'' a dream, but a direction worth moving in: not artificial churches, but new forms of community.

\responsehead

The post is a proposal, and it states its own uncertainty at each step. Its premise is a guess, hedged with ``probably''; its design ideas are things ``you could suspect'' that ``would be better tested experimentally''; its conclusion is rated ``Maybe. Probably.'' a dream. It should be judged as that.

The premise has support in general form. Lim and Putnam, writing in 2010, found that religious people's greater life satisfaction comes from attending services and having friends in their congregation, with little role for private belief. One sentence goes further than this: churches do not admit that community is ``nearly all of what people get out of it.'' Regular churchgoers, asked by Pew in 2018, mostly named becoming closer to God as their main reason, and about one in twenty named community. Stated reasons and measured benefits are different things, so neither finding settles the matter, but ``nearly all'' goes beyond both.

The definition of ``optimal'' does useful work. The list of church practices is judged only as a design for community, and the post grants that spending on religion makes sense for those who want religion. For information, one proposal is already common practice: studies by Partners for Sacred Places, an organization that supports historic religious buildings, report that most people served by programs in church buildings are not members, and that in one study only 9 per cent of visits were for worship.

The turn from community for its own sake to communities around a shared goal is argued in the text, which corrects itself when it notices that hunter-gatherer bands had a purpose. The same research bears on it. Lim and Putnam found that the benefit of friendship in a congregation depended on a strong religious identity. Whether a shared secular goal can supply that identity is the open question for the taskforce proposal, and the post does not claim to know.

In short: a hedged proposal for post-religious communities built around shared work, whose premise, that community is what churches chiefly give, has support in research on religion and life satisfaction.
\end{afterword}
